\scanpage{240}
\mspage{190}

\diagram{\begingroup\scriptsize
\begin{tikzcd}[ampersand replacement=\&,column sep=.65em,row sep=.35em]
0\arrow[r,no head] \& H^{n-1}(X,M)\arrow[r,no head]\arrow[d,equal] \& H^{n-1}(E,M)\arrow[r,no head]\arrow[d,equal] \& H^0(X,H^{n-1}(Z,M))\arrow[r] \& H^n(X,M)\arrow[r]\arrow[d,equal] \& H^n(E\unsure{,|),}M
\\
\&0\&0\&\&0\&
\end{tikzcd}\endgroup}

\[
\text{donc}\qquad H^0(X,\underbrace{H^{n-1}(Z,M)}_{\substack{\text{système local}\\\text{\uline{constant} sur }X\\\text{car }X\text{ est }1\text{-connexe}}})=0
\]

\[
\text{\unsure{donc}}\qquad H^{n-1}(Z,M)=0
\]

ce qui, par hyp. de récurrence, implique $\pi_{n-1}(Z,\unsure{x})=0$, donc $\pi_n(X,x)=0$, qed.

\uline{Corollaire} Pour que $X$ soit $W_\infty$-asphérique, \unsure{i.e.} pour qu'on ait \deleted{\unsure{pour}} $G\xrightarrow{\sim}H^0(X,G)$, $H^1(\unsure{X},G)=0$ tout groupe, et $H^i(X,G)=0$ si $i\geqslant 2$ pour $G$ abélien, il f. et s. que \deleted{\unsure{les}}

$1^\circ)$ $X$ est $\unsure{0{=}|0\text{-}}$connexe, et ($x$ étant un point de $X$) $2^\circ)$ les $\Omega^i(X,x)$ $0$-connexes pour $i\geqslant 1$.
% archive sheet number: 240
